% !TeX program = xelatex
% 2025 CUMCM A paper, electronic edition: abstract is page 1.
\documentclass[UTF8,zihao=-4]{ctexart}

\usepackage[a4paper,top=2.5cm,bottom=2.5cm,left=2.5cm,right=2.5cm]{geometry}
\usepackage{amsmath,amssymb,bm}
\usepackage{booktabs,array,longtable,tabularx}
\usepackage[table]{xcolor}
\usepackage{graphicx}
\usepackage{float}
\usepackage{fancyhdr}
\usepackage{enumitem}
\usepackage{listings}
\usepackage{pdflscape}
\usepackage[hidelinks]{hyperref}
\usepackage{microtype}

\graphicspath{{../outputs/figures/}}
\setlength{\parindent}{2em}
\setlength{\parskip}{0pt}
\linespread{1.25}
\setlength{\emergencystretch}{2em}
\setlist{nosep,leftmargin=2.5em}
\renewcommand{\arraystretch}{1.16}
\setcounter{secnumdepth}{3}

\pagestyle{fancy}
\fancyhf{}
\fancyfoot[C]{\thepage}
\renewcommand{\headrulewidth}{0pt}

\ctexset{
  section={format=\centering\heiti\zihao{3},beforeskip=1.2ex,afterskip=0.8ex},
  subsection={format=\heiti\zihao{4},beforeskip=1.0ex,afterskip=0.5ex}
}

\lstdefinestyle{python}{
  language=Python,
  basicstyle=\ttfamily\zihao{6},
  numbers=left,
  numberstyle=\tiny,
  numbersep=6pt,
  frame=single,
  breaklines=true,
  breakatwhitespace=false,
  showstringspaces=false,
  columns=fullflexible,
  keepspaces=true,
  tabsize=4
}

\newcommand{\PaperTitle}[1]{%
  \begin{center}
    {\heiti\zihao{2} #1}\par
  \end{center}
  \vspace{0.8em}
  \noindent\textbf{摘要：}%
}

\newif\iffullsourceappendix
\fullsourceappendixtrue

\begin{document}
\pagenumbering{arabic}
\setcounter{page}{1}

\PaperTitle{烟幕干扰弹投放策略的时空协同优化}

针对无人机投放烟幕干扰弹以遮蔽圆柱形真目标的问题，本文建立导弹、无人机、烟幕弹与烟幕云团的统一三维运动学模型，并以“来袭导弹到真目标的全部视线均被烟幕阻断”作为有效遮蔽判据。区别于只考察目标中心视线的近似方法，本文将真目标视为具有实际尺寸的圆柱体：单云团情形利用烟幕球对导弹视点形成的切锥，将整圆柱遮蔽转化为上、下底圆周上的连续极值问题；多云团情形采用“对每条目标视线，至少存在一个有效云团与之相交”的联合覆盖口径，即 $\forall\text{-ray}\ \exists\text{-cloud}$，从而允许不同云团共同覆盖目标视域。有效遮蔽时长由连续时间区间的并集测度计算，重叠时段不重复计时，不连续时段可以累加。

对于问题一，在 $g=9.8\ \mathrm{m/s^2}$ 下求得投放点为 $(17620,0,1800)$ m，起爆点为 $(17188,0,1736.496)$ m，严格整圆柱遮蔽区间为 $[8.056445,\allowbreak{}9.448088]$ s，有效时长为 $1.391643$ s。目标中心视线近似给出 $1.435082$ s，比严格结果多计约 $3.121\%$。对于问题二，采用中心视线代理筛选、差分进化广域搜索、多起点无梯度局部精修和严格几何统一复算，得到当前最佳已知时长 $4.542880$ s。对于问题三，在同机三弹及投弹间隔约束下，得到联合遮蔽区间 $[2.764945,\allowbreak{}8.278596]$ s，总时长 $5.513651$ s；当前候选实质由前两枚弹完成，不能据此断言第三枚弹在全局最优策略中无用。对于问题四，三架无人机各投一弹形成三段接力遮蔽，联合总时长为 $10.138857$ s。对于问题五，以三枚导弹遮蔽时长之和为主目标，通过单弹候选库、分配组合、逐机块精修和边际填充构成 40 维分层求解，得到 $(D_1,\allowbreak{}D_2,\allowbreak{}D_3)=(10.250038,\allowbreak{}1.912852,\allowbreak{}1.741134)$ s，$J_{\mathrm{sum}}=13.904024$ s，公平性诊断量 $J_{\min}=1.741134$ s；最终五架无人机各使用一枚有效弹。各问题均进行了可行性、网格收敛、边界探针、独立视线—球相交和内核回归检查。由于问题二至问题五的优化均受评估预算限制，本文只将其结果称为经过严格复算的最佳已知策略，不作全局最优声明。


\noindent\textbf{关键词：}烟幕干扰；视线遮蔽；圆柱目标；联合覆盖；区间并集；非光滑优化


\vfill
\noindent{\footnotesize\color{gray!75!black}
\textbf{AI 辅助标注：}本文第 1--13 节的题面整理、模型规格检查、代码实现、数值验证与文字初稿使用了 OpenAI Codex（GPT-5）辅助\cite{codex}；判据选择、参数口径和最终取舍由使用者确认，数值由随文程序复算。关键交互与人工修改情况见支撑材料《AI工具使用详情》。}

\clearpage
\input{body.tex}

\section*{参考文献}
\addcontentsline{toc}{section}{参考文献}
\begin{thebibliography}{9}
  \bibitem{cumcm2025a} 全国大学生数学建模竞赛组委会. 2025 年全国大学生数学建模竞赛 A 题：烟幕干扰弹的投放策略[Z]. 2025.
  \bibitem{codex} OpenAI Codex, GPT-5, OpenAI, 使用日期：2026-08-20--2026-08-21.
\end{thebibliography}

\appendix
\section{支撑材料文件列表}

本论文的支撑材料与论文电子版分开提交。主要文件如下：
\begin{longtable}{p{0.32\linewidth}p{0.60\linewidth}}
\toprule
文件或目录 & 内容说明 \\
\midrule
\texttt{problem.md} & 经页面核对的结构化题面 \\
\texttt{model-spec.md} & 模型边界、变量、约束、判据与需求追踪 \\
\texttt{configs/} & 问题一至问题五的正式运行配置 \\
\texttt{src/q1.py}--\texttt{src/q5.py} & 五个问题的完整可运行源程序 \\
\texttt{run.py} & 统一运行入口 \\
\texttt{tests/} & 联合覆盖、事件边界与回归测试 \\
\texttt{outputs/manifest.json} & 正式运行种子和产物清单 \\
\texttt{outputs/logs/} & 优化状态、边界、收敛和独立几何验证记录 \\
\texttt{outputs/tables/} & 结果表、区间表及三个结果工作簿 \\
\texttt{AI工具使用详情.pdf} & AI 工具名称、用途、关键交互与人工修改说明 \\
\bottomrule
\end{longtable}

\section{完整可运行源程序}

以下程序与支撑材料中的源文件一致。运行环境和命令见支撑材料中的
\texttt{README.md}。为满足论文附录包含完整程序的要求，正式提交稿默认展开全部源代码；若仅进行正文排版预览，可将导言区的
\verb|\fullsourceappendixtrue| 改为 \verb|\fullsourceappendixfalse|。

\iffullsourceappendix
  \subsection{统一运行入口}
  \lstinputlisting[style=python]{../run.py}
  \subsection{问题一程序}
  \lstinputlisting[style=python]{../src/q1.py}
  \subsection{问题二程序}
  \lstinputlisting[style=python]{../src/q2.py}
  \subsection{问题三程序}
  \lstinputlisting[style=python]{../src/q3.py}
  \subsection{问题四程序}
  \lstinputlisting[style=python]{../src/q4.py}
  \subsection{问题五程序}
  \lstinputlisting[style=python]{../src/q5.py}
\else
  本预览稿未展开源代码；正式提交前必须重新启用完整源程序附录。
\fi

\end{document}
